\subsection{Heuristic candidate sets have low corridor coverage}
\paragraph{GT is well aligned to the road graph.}
On Detroit segments, ground-truth points are typically close to the road graph: $94\%$ of GT points are within 1 grid cell of a road edge (p90 distance $<1$).
This suggests that low corridor match is not driven by GT/map misalignment.

\paragraph{$K$-shortest candidates miss many ground-truth corridors.}
Despite good GT alignment, heuristic candidate sets can have low corridor coverage.
With $K=20$ shortest-path variants, the median best-of-$K$ edge Jaccard match to GT is only $\approx 0.21$ on probed routes (Table~\ref{tab:candidate-coverage}), indicating that ``classify among $K$'' can be an unreliable surrogate for learning the corridor distribution.
This motivates learning structured corridor commitments rather than depending on a fixed candidate set.

\begin{table}[t]
  \centering
  \small
  \begin{tabular}{lccc}
    \toprule
    Setting & $K$ & $J_{\mathrm{best}}$ (p50) & $J_{\mathrm{best}}$ (p90) \\
    \midrule
    Detroit, $K$-shortest candidates & 20 & 0.215 & 0.699 \\
    \bottomrule
  \end{tabular}
  \caption{Candidate-set coverage diagnostic (Detroit). Even when GT points are close to the road graph, $K$-shortest candidates can have low best-of-$K$ edge Jaccard match to GT corridors.}
  \label{tab:candidate-coverage}
\end{table}

\subsection{Context carries measurable corridor-level signal (gate)}
\paragraph{AUC on corridor clusters.}
Using OD buckets and edge-Jaccard clustering within each bucket, we form a corridor-label prediction gate.
On Columbus, time+tier features achieve a strong signal (mean AUC $\approx 0.84$ over 18 usable OD groups).
Detroit exhibits a weaker but above-chance signal (mean AUC $\approx 0.61$ over 4 groups), highlighting the role of data volume and OD repetition.

\begin{table}[t]
  \centering
  \small
  \begin{tabular}{lccc}
    \toprule
    City & Usable OD groups & AUC(time+tier) mean & Gate \\
    \midrule
    Detroit & 4 & 0.613 & GO \\
    Columbus & 18 & 0.844 & STRONG\_GO \\
    \bottomrule
  \end{tabular}
  \caption{Semantic informativeness gate (cluster-based). AUC is computed for predicting the top-2 corridor clusters within multi-modal OD buckets using time and road-tier features.}
  \label{tab:semantic-gate}
\end{table}

\subsection{Graph-based corridor generation: node AR vs waypoint AR + A*}
\paragraph{Node-level AR is fragile at long horizons.}
We train a next-node predictor on graph transitions (val accuracy $\approx 0.76$), but long-horizon sampling remains challenging: for $K=20$ samples, best-of-$K$ corridor match is modest ($J_{\mathrm{best}}$ mean $\approx 0.20$) and success depends strongly on route length.
This reflects compounding error over hundreds of steps.

\paragraph{Waypoint AR + A* reduces the decision horizon but needs reachability-aware decoding.}
Our waypoint AR model predicts a 4-step waypoint commitment in a coarse $32\times 32$ bin space (best validation accuracy $\approx 0.34$).
Executing sampled waypoints with A* guarantees graph-validity for successful samples, but the current prototype has a low success rate because sampled waypoint bins can map to unreachable node pairs.
This indicates the bottleneck is reachability constraints rather than long-horizon exposure bias.

\begin{table}[t]
  \centering
  \small
  \begin{tabular}{lccc}
    \toprule
    Method & Success (mean) & $J_{\mathrm{best}}$ (mean) & Notes \\
    \midrule
    Node-level AR (T3) & 0.376 & 0.204 & long-horizon drift \\
    Waypoint AR + A* (T4) & 0.115 & 0.179 & unreachable waypoint segments \\
    \bottomrule
  \end{tabular}
  \caption{Graph corridor generation on the combined Detroit+Columbus dataset (200 sampled routes, $K=20$). Success denotes completion to $d$; $J_{\mathrm{best}}$ is best-of-$K$ edge Jaccard to the GT corridor.}
  \label{tab:t3-t4}
\end{table}
